\chapter{矩阵相似理论}

\section{逐题解析}

\subsection{5.1 特征值与特征向量}

\begin{solutionblock}
\textbf{5.1 例 1}

\question{求矩阵 \(A=\begin{pmatrix}1&2&-3\\-1&4&-3\\1&-2&5\end{pmatrix}\) 的全部特征值及其对应的特征向量。}



\analysis{设
\[
A=\begin{pmatrix}
1&2&-3\\
-1&4&-3\\
1&-2&5
\end{pmatrix}.
\]
计算特征多项式：
\[
|\lambda E-A|
=\begin{vmatrix}
\lambda-1&-2&3\\
1&\lambda-4&3\\
-1&2&\lambda-5
\end{vmatrix}
=(\lambda-6)(\lambda-2)^2.
\]
所以特征值为 \(\lambda=2,2,6\)。

当 \(\lambda=2\) 时，
\[
A-2E=\begin{pmatrix}
-1&2&-3\\
-1&2&-3\\
1&-2&3
\end{pmatrix},
\quad -x_1+2x_2-3x_3=0.
\]
令 \(x_2=s,x_3=t\)，得 \(x_1=2s-3t\)，故
\[
\xi=s(2,1,0)^T+t(-3,0,1)^T.
\]
当 \(\lambda=6\) 时，解 \((A-6E)x=0\)，可得
\[
\xi=k(-1,-1,1)^T.
\]}
\answer{\[
\lambda=2:\ k_1(2,1,0)^T+k_2(-3,0,1)^T,\qquad
\lambda=6:\ k(-1,-1,1)^T,
\]
其中对应特征向量要求系数不全为零。}
\examnote{二重特征值要继续求特征子空间维数，不能只写一个特征向量。}
\end{solutionblock}

\begin{solutionblock}
\textbf{5.1 例 2}

\question{分别求下列矩阵的特征值：\((1)\ \begin{pmatrix}1&0&0\\0&2&0\\6&-2&3\end{pmatrix}\)；\((2)\ \begin{pmatrix}1&2&3\\2&4&6\\3&6&9\end{pmatrix}\)。}



\analysis{第 (1) 个矩阵
\[
\begin{pmatrix}
1&0&0\\
0&2&0\\
6&-2&3
\end{pmatrix}
\]
是下三角矩阵，特征值就是主对角线元素：
\[
1,\ 2,\ 3.
\]

第 (2) 个矩阵
\[
\begin{pmatrix}
1&2&3\\
2&4&6\\
3&6&9
\end{pmatrix}
=\alpha\alpha^T,\qquad \alpha=(1,2,3)^T.
\]
秩一矩阵 \(\alpha\alpha^T\) 的非零特征值为
\[
\alpha^T\alpha=1^2+2^2+3^2=14,
\]
其余特征值为 \(0\)。}
\answer{第 (1) 个矩阵的特征值为 \(1,2,3\)；第 (2) 个矩阵的特征值为 \(14,0,0\)。}
\end{solutionblock}

\begin{solutionblock}
\textbf{5.1 例 3}

\question{已知 \(0\) 是矩阵 \(A=\begin{pmatrix}1&0&1\\0&2&0\\1&0&a\end{pmatrix}\) 的一个特征值，求 \(a\)，并求此时 \(A\) 的全部特征值。}



\analysis{\(0\) 是
\[
A=\begin{pmatrix}
1&0&1\\
0&2&0\\
1&0&a
\end{pmatrix}
\]
的特征值，等价于 \(|A|=0\)。计算
\[
|A|=2\begin{vmatrix}1&1\\1&a\end{vmatrix}=2(a-1),
\]
所以 \(a=1\)。此时
\[
A=\begin{pmatrix}
1&0&1\\
0&2&0\\
1&0&1
\end{pmatrix}.
\]
第一、三坐标构成的块为
\[
\begin{pmatrix}1&1\\1&1\end{pmatrix},
\]
特征值为 \(0,2\)，中间坐标给出特征值 \(2\)。}
\answer{\[
a=1,\qquad A\text{ 的全部特征值为 }0,2,2.
\]}
\end{solutionblock}

\begin{solutionblock}
\textbf{5.1 例 4}

\question{设矩阵 \(A\) 可逆，且 \(A^{-1}\) 的特征值为 \(1,-\frac12,\frac13\)，求 \(A\)、\((2A)^*\)、\(A^2+4A+E\) 的特征值。}



\analysis{\(A^{-1}\) 的特征值为 \(1,-\frac12,\frac13\)，所以 \(A\) 的特征值为其倒数：
\[
1,\ -2,\ 3.
\]

对于 \(A\)，若特征值为 \(\lambda_i\)，则 \(A^*\) 的对应特征值为
\[
\frac{|A|}{\lambda_i}.
\]
这里
\[
|A|=1\cdot(-2)\cdot3=-6,
\]
故 \(A^*\) 的特征值为
\[
-6,\ 3,\ -2.
\]
三阶矩阵满足
\[
(2A)^*=2^{3-1}A^*=4A^*,
\]
所以 \((2A)^*\) 的特征值为
\[
-24,\ 12,\ -8.
\]

对多项式矩阵 \(A^2+4A+E\)，特征值由函数
\[
f(x)=x^2+4x+1
\]
作用在 \(A\) 的特征值上得到：
\[
f(1)=6,\qquad f(-2)=-3,\qquad f(3)=22.
\]}
\answer{\[
A:\ 1,-2,3;\qquad
(2A)^*:\ -24,12,-8;\qquad
A^2+4A+E:\ 6,-3,22.
\]}
\end{solutionblock}

\begin{solutionblock}
\textbf{5.1 例 5}

\question{设三阶矩阵 \(A\) 满足 \(|A-E|=|A+2E|=|2A+3E|=0\)，求 \(|2A^*-3E|\)。}



\analysis{由
\[
|A-E|=0,\quad |A+2E|=0,\quad |2A+3E|=0
\]
知 \(A\) 的三个特征值分别为
\[
1,\ -2,\ -\frac32.
\]
故
\[
|A|=1\cdot(-2)\cdot\left(-\frac32\right)=3.
\]
因为 \(A\) 可逆，\(A^*\) 的特征值为 \(|A|/\lambda_i\)，即
\[
3,\ -\frac32,\ -2.
\]
于是 \(2A^*-3E\) 的特征值为
\[
2\cdot3-3=3,\qquad
2\cdot\left(-\frac32\right)-3=-6,\qquad
2\cdot(-2)-3=-7.
\]
所求行列式为三个特征值之积：
\[
|2A^*-3E|=3\cdot(-6)\cdot(-7)=126.
\]}
\answer{\(126\)}
\end{solutionblock}

\begin{solutionblock}
\textbf{5.1 例 6}

\question{设 \(A\) 为 \(n\) 阶可逆矩阵，\(A\) 的每行元素之和均为 \(a\,(a\ne0)\)，且 \(|A|=1\)。求 \(A^*\) 的每行元素之和。}



\analysis{设 \(\mathbf{1}=(1,1,\ldots,1)^T\)。每行元素之和均为 \(a\)，等价于
\[
A\mathbf{1}=a\mathbf{1}.
\]
由于 \(A\) 可逆且 \(a\ne0\)，两边左乘 \(A^{-1}\)，得
\[
A^{-1}\mathbf{1}=\frac1a\mathbf{1}.
\]
又 \(|A|=1\)，所以
\[
A^*=|A|A^{-1}=A^{-1}.
\]
矩阵乘以 \(\mathbf{1}\) 得到的正是各行元素之和，因此 \(A^*\) 的每行元素之和均为 \(\frac1a\)。}
\answer{\(\dfrac1a\)}
\end{solutionblock}

\begin{solutionblock}
\textbf{5.1 例 7}

\question{设
\[
A=\begin{pmatrix}
a&-1&c\\
5&b&3\\
1-c&0&-a
\end{pmatrix},
\qquad |A|=-1.
\]
若 \(A\) 的伴随矩阵 \(A^*\) 有特征值 \(\lambda_0\)，且属于 \(\lambda_0\) 的特征向量为 \(\alpha=(-1,-1,1)^T\)，求 \(a,b,c\) 及 \(\lambda_0\)。}



\analysis{设
\[
A=\begin{pmatrix}
a&-1&c\\
5&b&3\\
1-c&0&-a
\end{pmatrix},\qquad \alpha=(-1,-1,1)^T.
\]
因为 \(\alpha\) 是 \(A^*\) 属于 \(\lambda_0\) 的特征向量，且 \(|A|=-1\ne0\)，所以 \(\alpha\) 也是 \(A\) 的特征向量。设
\[
A\alpha=\mu\alpha.
\]
直接计算
\[
A\alpha=
\begin{pmatrix}
-a+1+c\\
-5-b+3\\
-1+c-a
\end{pmatrix}
=
\begin{pmatrix}
-a+c+1\\
-b-2\\
c-a-1
\end{pmatrix}.
\]
与 \(\mu(-1,-1,1)^T\) 比较，得
\[
-a+c+1=-\mu,\qquad -b-2=-\mu,\qquad c-a-1=\mu.
\]
第一式与第三式给出
\[
-a+c+1=-(c-a-1),
\]
恒成立；再由第二、第三式得
\[
\mu=c-a-1=b+2.
\]
即
\[
b=c-a-3.
\]
再代入 \(|A|=-1\)，可解得
\[
a=2,\qquad b=-3,\qquad c=2,\qquad \mu=-1.
\]
最后利用 \(A^*=|A|A^{-1}\)：若 \(A\alpha=\mu\alpha\)，则
\[
A^*\alpha=\frac{|A|}{\mu}\alpha=\frac{-1}{-1}\alpha=\alpha.
\]
所以 \(\lambda_0=1\)。}
\answer{\[
a=2,\qquad b=-3,\qquad c=2,\qquad \lambda_0=1.
\]}
\examnote{可逆时 \(A\) 与 \(A^*\) 的特征向量相同，特征值按 \(\lambda\mapsto |A|/\lambda\) 转换。}
\end{solutionblock}

\begin{solutionblock}
\textbf{5.1 例 8}

\question{设 \(A\) 为 \(n\) 阶实对称矩阵，\(P\) 为 \(n\) 阶可逆矩阵。已知 \(n\) 维列向量 \(\alpha\) 是矩阵 \(A\) 属于特征值 \(\lambda\) 的特征向量，判断矩阵 \((P^{-1}AP)^T\) 属于特征值 \(\lambda\) 的特征向量：
\[
\text{A. }P^{-1}\alpha\qquad
\text{B. }P^T\alpha\qquad
\text{C. }P\alpha\qquad
\text{D. }(P^{-1})^T\alpha.
\]}



\analysis{设 \(A\alpha=\lambda\alpha\)。因为 \(A\) 为实对称矩阵，故 \(A^T=A\)。记
\[
C=(P^{-1}AP)^T=P^TA^T(P^{-1})^T=P^TA(P^{-1})^T.
\]
考察 \(P^T\alpha\)：
\[
C(P^T\alpha)
=P^TA(P^{-1})^TP^T\alpha
=P^TA\alpha
=\lambda P^T\alpha.
\]
故 \(P^T\alpha\) 是 \(C\) 属于 \(\lambda\) 的特征向量。}
\answer{选 B，\(P^T\alpha\)。}
\end{solutionblock}

\begin{solutionblock}
\textbf{5.1 例 9}

\question{若 \(\alpha_1\) 是矩阵 \(A\) 属于特征值 \(\lambda=1\) 的特征向量，\(\alpha_2,\alpha_3\) 是矩阵 \(A\) 属于特征值 \(\lambda=2\) 的两个线性无关特征向量。设
\[
AP=P\begin{pmatrix}
1&0&0\\
0&2&0\\
0&0&2
\end{pmatrix},
\]
判断矩阵 \(P\) 不可以是哪一个：
\[
\text{A. }(\alpha_1,-2\alpha_2,\alpha_3),\quad
\text{B. }(\alpha_1,\alpha_2+\alpha_3,\alpha_2-\alpha_3),
\]
\[
\text{C. }(\alpha_1,\alpha_3,\alpha_2),\quad
\text{D. }(\alpha_1+\alpha_2,\alpha_1-\alpha_2,\alpha_3).
\]}



\analysis{由
\[
AP=P\begin{pmatrix}
1&0&0\\
0&2&0\\
0&0&2
\end{pmatrix}
\]
可知 \(P\) 的第一列必须是属于 \(\lambda=1\) 的特征向量，第二、三列必须是属于 \(\lambda=2\) 的两个线性无关特征向量。

选项 A、B、C 中第一列仍为 \(\alpha_1\)，第二、三列均是 \(\alpha_2,\alpha_3\) 的线性无关组合，所以可以作为 \(P\)。

选项 D 第一列为 \(\alpha_1+\alpha_2\)。由于
\[
A(\alpha_1+\alpha_2)=\alpha_1+2\alpha_2
\]
不等于 \(1(\alpha_1+\alpha_2)\)，也不等于 \(2(\alpha_1+\alpha_2)\)，故它不是单一特征值对应的特征向量，不能作为第一列。}
\answer{选 D。}
\end{solutionblock}

\subsection{5.2 秩为 1 矩阵专题}

\begin{solutionblock}
\textbf{5.2 例 1}

\question{设 \(\alpha=(1,2,3)^T\)，\(\beta=(1,\frac12,\frac13)^T\)，\(A=\alpha\beta^T\)。求 \(A^n\)。}



\analysis{设
\[
\alpha=(1,2,3)^T,\qquad \beta=\left(1,\frac12,\frac13\right)^T,\qquad A=\alpha\beta^T.
\]
秩一矩阵满足
\[
A^2=(\alpha\beta^T)(\alpha\beta^T)
=\alpha(\beta^T\alpha)\beta^T
=(\beta^T\alpha)A.
\]
这里
\[
\beta^T\alpha=1+1+1=3.
\]
因此
\[
A^n=(\beta^T\alpha)^{n-1}A=3^{n-1}A,\qquad n\ge1.
\]}
\answer{\[
A^n=3^{n-1}\alpha\beta^T=3^{n-1}A.
\]}
\end{solutionblock}

\begin{solutionblock}
\textbf{5.2 例 2}

\question{设 \(A=E+\alpha\beta^T\)，且 \(\alpha^T\beta=2\)。求 \(A^{-1}\)。}



\analysis{令 \(A=E+\alpha\beta^T\)。秩一修正公式为
\[
(E+\alpha\beta^T)^{-1}
=E-\frac{\alpha\beta^T}{1+\beta^T\alpha},
\]
前提是 \(1+\beta^T\alpha\ne0\)。题中
\[
\alpha^T\beta=\beta^T\alpha=2,
\]
故
\[
1+\beta^T\alpha=3\ne0,
\]
于是 \(A\) 可逆，且
\[
A^{-1}=E-\frac13\alpha\beta^T.
\]}
\answer{\[
A^{-1}=E-\frac13\alpha\beta^T.
\]}
\examnote{考研中遇到 \(E+\alpha\beta^T\)，优先想到秩一修正公式；也可设逆为 \(E+k\alpha\beta^T\) 后比较系数求 \(k\)。}
\end{solutionblock}

\begin{solutionblock}
\textbf{5.2 例 3}

\question{设 \(\alpha\beta^T\) 为三阶矩阵，且 \(\alpha\beta^T\) 与
\[
\begin{pmatrix}2&0&0\\0&0&0\\0&0&0\end{pmatrix}
\]
相似。求 \(\beta^T\alpha\)。}



\analysis{\(\alpha\beta^T\) 为三阶秩一矩阵，非零特征值为
\[
\beta^T\alpha.
\]
它与
\[
\begin{pmatrix}
2&0&0\\
0&0&0\\
0&0&0
\end{pmatrix}
\]
相似，故两者特征值相同，为 \(2,0,0\)。因此
\[
\beta^T\alpha=2.
\]}
\answer{\(2\)}
\end{solutionblock}

\begin{solutionblock}
\textbf{5.2 例 4}

\question{设
\[
A=\begin{pmatrix}2&1&3\\4&2&6\\6&3&9\end{pmatrix},
\qquad
B=A+E.
\]
求 \(A\) 与 \(B\) 的特征值和对应特征向量。}



\analysis{设
\[
A=\begin{pmatrix}2&1&3\\4&2&6\\6&3&9\end{pmatrix}
=\alpha\beta^T,\quad
\alpha=(1,2,3)^T,\quad \beta=(2,1,3)^T.
\]
则
\[
\beta^T\alpha=2+2+9=13.
\]
所以 \(A\) 的特征值为 \(13,0,0\)。对应 \(\lambda=13\) 的特征向量可取 \(\alpha=(1,2,3)^T\)；对应 \(\lambda=0\) 的特征向量满足
\[
\beta^Tx=2x_1+x_2+3x_3=0,
\]
即
\[
x=s(-1,2,0)^T+t(-3,0,2)^T.
\]

对
\[
B=\begin{pmatrix}3&1&3\\4&3&6\\6&3&10\end{pmatrix}
=A+E.
\]
若 \(A\) 的特征值为 \(\lambda\)，则 \(A+E\) 的特征值为 \(\lambda+1\)。因此 \(B\) 的特征值为
\[
14,1,1.
\]
特征向量不变：\(\lambda=14\) 对应 \((1,2,3)^T\)；\(\lambda=1\) 对应
\[
2x_1+x_2+3x_3=0.
\]}
\answer{\[
A:\ \lambda=13,\ \xi=k(1,2,3)^T;\quad
\lambda=0,\ \xi=s(-1,2,0)^T+t(-3,0,2)^T.
\]
\[
B:\ \lambda=14,\ \xi=k(1,2,3)^T;\quad
\lambda=1,\ \xi=s(-1,2,0)^T+t(-3,0,2)^T.
\]}
\end{solutionblock}

\subsection{5.3 矩阵相似}

\begin{solutionblock}
\textbf{5.3 例 1}

\question{已知矩阵
\[
A=\begin{pmatrix}-2&-2&1\\2&x&-2\\0&0&-2\end{pmatrix},
\qquad
B=\begin{pmatrix}2&1&0\\0&-1&0\\0&0&y\end{pmatrix}
\]
相似，求 \(x,y\)。}



\analysis{相似矩阵有相同的特征多项式。对
\[
A=\begin{pmatrix}
-2&-2&1\\
2&x&-2\\
0&0&-2
\end{pmatrix},
\quad
B=\begin{pmatrix}
2&1&0\\
0&-1&0\\
0&0&y
\end{pmatrix},
\]
分别计算
\[
|\lambda E-A|
=-(\lambda+2)(-\lambda^2+\lambda x-2\lambda+2x-4),
\]
\[
|\lambda E-B|=(\lambda-2)(\lambda+1)(\lambda-y).
\]
比较系数可得
\[
x=3,\qquad y=-2.
\]
此时两个矩阵的特征值为 \(2,-1,-2\)，互异，故均可对角化并相似于同一个对角矩阵。}
\answer{\[
x=3,\qquad y=-2.
\]}
\end{solutionblock}

\begin{solutionblock}
\textbf{5.3 例 2}

\question{设
\[
B=\begin{pmatrix}0&0&1\\0&1&0\\1&0&0\end{pmatrix},
\qquad A\sim B.
\]
求 \(r(A-2E)+r(A-E)\)。}



\analysis{题设 \(A\sim B\)，其中
\[
B=\begin{pmatrix}
0&0&1\\
0&1&0\\
1&0&0
\end{pmatrix}.
\]
相似变换不改变矩阵 \(A-\lambda E\) 的秩，因为
\[
A=P^{-1}BP
\quad\Longrightarrow\quad
A-\lambda E=P^{-1}(B-\lambda E)P.
\]
矩阵 \(B\) 的特征值为 \(1,1,-1\)。所以 \(2\) 不是特征值，
\[
r(A-2E)=r(B-2E)=3.
\]
而 \(\lambda=1\) 的几何重数为 \(2\)，故
\[
\dim\ker(B-E)=2,\qquad r(B-E)=3-2=1.
\]
于是
\[
r(A-2E)+r(A-E)=3+1=4.
\]}
\answer{选 C，\(4\)。}
\end{solutionblock}

\begin{solutionblock}
\textbf{5.3 例 3}

\question{设 \(A=\alpha\alpha^T\)，其中 \(\alpha=(1,2,-1)^T\)，且 \(A\sim B\)。求 \(|(B+E)^*|\)。}



\analysis{已知
\[
\alpha=(1,2,-1)^T,\qquad A=\alpha\alpha^T.
\]
秩一矩阵 \(A\) 的非零特征值为
\[
\alpha^T\alpha=1+4+1=6,
\]
故 \(A\) 的特征值为 \(6,0,0\)。由于 \(A\sim B\)，\(B\) 的特征值也为 \(6,0,0\)，从而 \(B+E\) 的特征值为
\[
7,1,1.
\]
于是
\[
|B+E|=7.
\]
三阶矩阵 \(C\) 满足
\[
|C^*|=|C|^{3-1}=|C|^2.
\]
取 \(C=B+E\)，得
\[
|(B+E)^*|=7^2=49.
\]}
\answer{\(49\)}
\end{solutionblock}

\begin{solutionblock}
\textbf{5.3 例 4}

\question{若 \(A\) 与 \(B\) 相似，判断有关特征矩阵、可逆性、特征向量和可对角化性的命题中哪一个一定正确。}



\analysis{若 \(A\sim B\)，则 \(A=P^{-1}BP\)。因此
\[
A-\lambda E=P^{-1}(B-\lambda E)P.
\]
取 \(\lambda=0\) 可知 \(A\) 与 \(B\) 同时可逆或同时不可逆，所以选项 B 正确。

选项 A 写成矩阵相等，过强；相似只能推出 \(\lambda E-A\sim \lambda E-B\)，不能推出二者相等。选项 C 错在特征向量依赖具体坐标基，通常不会相同。选项 D 错在相似矩阵不一定可对角化。}
\answer{选 B。}
\end{solutionblock}

\begin{solutionblock}
\textbf{5.3 例 5}

\question{设 \(A,B\) 均可逆且 \(A\sim B\)，判断 \(A^T\sim B^T\)、\(A+A^T\sim B+B^T\)、\(A^{-1}\sim B^{-1}\)、\(A+A^{-1}\sim B+B^{-1}\) 中哪些成立。}



\analysis{设 \(A=P^{-1}BP\)。则
\[
A^T=P^TB^T(P^{-1})^T,
\]
所以 \(A^T\sim B^T\)，选项 A 正确。

若 \(A,B\) 可逆，则
\[
A^{-1}=P^{-1}B^{-1}P,
\]
所以 \(A^{-1}\sim B^{-1}\)，选项 C 正确。进一步
\[
A+A^{-1}=P^{-1}(B+B^{-1})P,
\]
故选项 D 正确。

但 \(A+A^T\) 中同时出现 \(P^{-1}BP\) 与 \(P^TB^T(P^{-1})^T\)，一般不能合并成同一个相似变换，所以不一定与 \(B+B^T\) 相似。}
\answer{选 B。}
\examnote{“多项式保持相似”要求表达式只由 \(A\) 的加减乘幂组成；一旦混入 \(A^T\)，就要格外小心。}
\end{solutionblock}

\begin{solutionblock}
\textbf{5.3 例 6}

\question{判断矩阵 \(\begin{pmatrix}1&1&0\\0&1&1\\0&0&1\end{pmatrix}\) 与给定四个三阶矩阵中的哪一个相似。}



\analysis{目标矩阵
\[
J=\begin{pmatrix}
1&1&0\\
0&1&1\\
0&0&1
\end{pmatrix}
\]
是三阶 Jordan 块。判断相似时，看
\[
N=J-E
\]
的秩结构：
\[
r(N)=2,\qquad r(N^2)=1.
\]
四个选项分别计算 \(M-E\) 与 \((M-E)^2\) 的秩，只有选项 A 满足
\[
r(M-E)=2,\qquad r((M-E)^2)=1.
\]
其余选项的 \(M-E\) 秩为 \(1\) 或平方为零，对应 Jordan 块结构不同。}
\answer{选 A。}
\end{solutionblock}

\begin{solutionblock}
\textbf{5.3 例 7}

\question{判断下列条件中哪一个是 \(A\sim B\) 的充分条件：\(A^*\sim B^*\)、\(A^T\sim B^T\)、\(f(A)\sim f(B)\)、\(AB\sim BA\)。}



\analysis{若 \(A\sim B\)，则必有 \(A^T\sim B^T\)；反过来，若 \(A^T\sim B^T\)，即存在可逆矩阵 \(P\) 使
\[
A^T=P^{-1}B^TP,
\]
两边转置得
\[
A=P^TB(P^{-1})^T,
\]
所以 \(A\sim B\)。因此 \(A^T\sim B^T\) 是 \(A\sim B\) 的充分条件。

选项 A 中伴随矩阵相似不能反推原矩阵相似；选项 C 中 \(f(A)\sim f(B)\) 信息可能丢失；选项 D 中 \(AB\sim BA\) 即使成立，也不能推出 \(A\sim B\)。}
\answer{选 B。}
\end{solutionblock}

\subsection{5.4 矩阵相似对角化}

\begin{solutionblock}
\textbf{5.4 例 1}

\question{设三阶矩阵 \(A\) 的特征值为 \(0,1,2\)，令 \(B=A^3-2A^2\)，求 \(r(B)\)。}



\analysis{设 \(f(x)=x^3-2x^2=x^2(x-2)\)，则
\[
B=A^3-2A^2=f(A).
\]
已知 \(A\) 的特征值为 \(0,1,2\)，且三个特征值互异，所以 \(A\) 可对角化。于是 \(B=f(A)\) 的特征值为
\[
f(0)=0,\qquad f(1)=-1,\qquad f(2)=0.
\]
\(B\) 只有一个非零特征值，故
\[
r(B)=1.
\]}
\answer{\(1\)}
\end{solutionblock}

\begin{solutionblock}
\textbf{5.4 例 2}

\question{设 \(\alpha=(1,1,1)^T\)，\(\beta=(1,0,k)^T\)，且 \(\alpha\beta^T\) 与 \(\operatorname{diag}(4,0,0)\) 相似，求 \(k\)。}



\analysis{秩一矩阵 \(\alpha\beta^T\) 的非零特征值为
\[
\beta^T\alpha.
\]
题设
\[
\alpha=(1,1,1)^T,\qquad \beta=(1,0,k)^T,
\]
所以
\[
\beta^T\alpha=1+0+k=1+k.
\]
又 \(\alpha\beta^T\) 与 \(\operatorname{diag}(4,0,0)\) 相似，故其非零特征值为 \(4\)。于是
\[
1+k=4,\qquad k=3.
\]}
\answer{\(3\)}
\end{solutionblock}

\begin{solutionblock}
\textbf{5.4 例 3}

\question{判断给定四个矩阵中哪一个不能相似对角化。}



\analysis{实对称矩阵一定可以正交相似对角化，所以选项 A 可对角化。

选项 B 为三角矩阵，特征值 \(1,2,3\) 互异，故可对角化。

选项 C 是秩一矩阵，特征值为 \(7,0,0\)，且 \(\lambda=0\) 的特征子空间维数为 \(2\)，也可对角化。

选项 D 的特征多项式为
\[
\lambda^2(\lambda-1),
\]
但
\[
r(D-0E)=2,\qquad \dim\ker D=1,
\]
\(\lambda=0\) 的几何重数小于代数重数 \(2\)，故不能对角化。}
\answer{选 D。}
\end{solutionblock}

\begin{solutionblock}
\textbf{5.4 例 4}

\question{设 \(A=\begin{pmatrix}4&6&0\\-3&-5&0\\-3&-6&1\end{pmatrix}\)，求可逆矩阵 \(P\)，使 \(P^{-1}AP\) 为对角矩阵。}



\analysis{设
\[
A=\begin{pmatrix}
4&6&0\\
-3&-5&0\\
-3&-6&1
\end{pmatrix}.
\]
计算
\[
|\lambda E-A|=(\lambda+2)(\lambda-1)^2.
\]
当 \(\lambda=-2\) 时，特征向量可取
\[
\xi_1=(-1,1,1)^T.
\]
当 \(\lambda=1\) 时，解 \((A-E)x=0\)，可取两个线性无关特征向量
\[
\xi_2=(-2,1,0)^T,\qquad \xi_3=(0,0,1)^T.
\]
令
\[
P=(\xi_1,\xi_2,\xi_3)
=\begin{pmatrix}
-1&-2&0\\
1&1&0\\
1&0&1
\end{pmatrix}.
\]
则 \(|P|\ne0\)，并且
\[
P^{-1}AP=\begin{pmatrix}
-2&0&0\\
0&1&0\\
0&0&1
\end{pmatrix}.
\]}
\answer{可取
\[
P=\begin{pmatrix}
-1&-2&0\\
1&1&0\\
1&0&1
\end{pmatrix},
\quad
P^{-1}AP=\operatorname{diag}(-2,1,1).
\]}
\end{solutionblock}

\begin{solutionblock}
\textbf{5.4 例 5}

\question{设 \(A=\begin{pmatrix}1&-1&1\\x&4&y\\-3&-3&5\end{pmatrix}\)。已知 \(A\) 有三个线性无关特征向量，且 \(\lambda=2\) 为二重特征值，求 \(x,y\)，并求 \(P\) 使 \(P^{-1}AP\) 为对角矩阵。}



\analysis{设
\[
A=\begin{pmatrix}
1&-1&1\\
x&4&y\\
-3&-3&5
\end{pmatrix}.
\]
由题意，\(A\) 有三个线性无关的特征向量，且 \(\lambda=2\) 是二重特征值，因此特征多项式应为
\[
(\lambda-2)^2(\lambda-\lambda_3).
\]
直接计算可得
\[
|\lambda E-A|=(\lambda-2)(\lambda^2-8\lambda+x+3y+16).
\]
若 \(\lambda=2\) 为二重根，则
\[
4-16+x+3y+16=0,\qquad x+3y+4=0. \tag{1}
\]
但仅有二重根还不够，还要求 \(\lambda=2\) 的特征子空间维数为 \(2\)。由
\[
A-2E=\begin{pmatrix}
-1&-1&1\\
x&2&y\\
-3&-3&3
\end{pmatrix}
\]
可知第一行与第三行成比例，要使秩为 \(1\)，第二行也必须与第一行成比例。结合 (1) 得
\[
x=2,\qquad y=-2.
\]
此时另一个特征值由迹确定：
\[
1+4+5=10=2+2+\lambda_3,\qquad \lambda_3=6.
\]
对应 \(\lambda=2\) 的特征向量可取
\[
\xi_1=(-1,1,0)^T,\qquad \xi_2=(1,0,1)^T,
\]
对应 \(\lambda=6\) 的特征向量可取
\[
\xi_3=(1,-2,3)^T.
\]
令
\[
P=(\xi_1,\xi_2,\xi_3)
=\begin{pmatrix}
-1&1&1\\
1&0&-2\\
0&1&3
\end{pmatrix}.
\]
则
\[
P^{-1}AP=\operatorname{diag}(2,2,6).
\]}
\answer{\[
x=2,\qquad y=-2,\qquad
P=\begin{pmatrix}
-1&1&1\\
1&0&-2\\
0&1&3
\end{pmatrix}.
\]}
\examnote{本题的关键陷阱是：二重特征值不代表一定能对角化，还要检查该特征值对应的线性无关特征向量个数。}
\end{solutionblock}

\begin{solutionblock}
\textbf{5.4 例 6}

\question{设 \(A=\begin{pmatrix}1&2&-3\\-1&4&-3\\1&a&5\end{pmatrix}\) 有二重特征值，求 \(a\)，并判断 \(A\) 是否可相似对角化。}



\analysis{设
\[
A=\begin{pmatrix}
1&2&-3\\
-1&4&-3\\
1&a&5
\end{pmatrix}.
\]
计算特征多项式：
\[
|\lambda E-A|=(\lambda-2)(\lambda^2-8\lambda+3a+18).
\]
若有一个二重根，则判别式为零或 \(\lambda=2\) 也是二次因子的根。等价计算可得
\[
a=-2\quad\text{或}\quad a=-\frac23.
\]

当 \(a=-2\) 时，
\[
|\lambda E-A|=(\lambda-6)(\lambda-2)^2,
\]
且 \(\lambda=2\) 的特征子空间维数为 \(2\)，故 \(A\) 可相似对角化。

当 \(a=-\frac23\) 时，
\[
|\lambda E-A|=(\lambda-4)^2(\lambda-2),
\]
但 \(\lambda=4\) 的特征子空间维数为 \(1\)，小于其代数重数 \(2\)，故不能相似对角化。}
\answer{\[
a=-2\ \text{或}\ -\frac23.
\]
其中 \(a=-2\) 时 \(A\) 可相似对角化；\(a=-\frac23\) 时不可相似对角化。}
\end{solutionblock}

\begin{solutionblock}
\textbf{5.4 例 7}

\question{设三阶矩阵 \(A\) 对应线性方程组 \(Ax=b\) 的通解为 \(k_1(-1,1,0)^T+k_2(2,0,1)^T+(1,1,-2)^T\)，其中 \(b=(6,6,-12)^T\)，求 \(A\)。}



\analysis{方程组 \(Ax=b\) 的通解为
\[
x=k_1(-1,1,0)^T+k_2(2,0,1)^T+(1,1,-2)^T.
\]
因此齐次方程组 \(Ax=0\) 的基础解系为
\[
\eta_1=(-1,1,0)^T,\qquad \eta_2=(2,0,1)^T.
\]
又特解 \(x_0=(1,1,-2)^T\) 满足
\[
Ax_0=b=(6,6,-12)^T=6x_0.
\]
所以 \(A\) 在基
\[
P=(\eta_1,\eta_2,x_0)
=\begin{pmatrix}
-1&2&1\\
1&0&1\\
0&1&-2
\end{pmatrix}
\]
下对应对角矩阵
\[
\operatorname{diag}(0,0,6).
\]
即
\[
AP=P\operatorname{diag}(0,0,6),
\quad
A=P\operatorname{diag}(0,0,6)P^{-1}.
\]
计算得
\[
A=\begin{pmatrix}
2&2&-2\\
2&2&-2\\
-4&-4&4
\end{pmatrix}.
\]}
\answer{\[
A=\begin{pmatrix}
2&2&-2\\
2&2&-2\\
-4&-4&4
\end{pmatrix}.
\]}
\end{solutionblock}

\begin{solutionblock}
\textbf{5.4 例 8}

\question{设 \(A=\begin{pmatrix}1&4&-2\\0&-1&0\\1&2&-2\end{pmatrix}\)，求 \(A^{100}\)。}



\analysis{设
\[
A=\begin{pmatrix}
1&4&-2\\
0&-1&0\\
1&2&-2
\end{pmatrix}.
\]
其特征值为 \(0,-1,-1\)，且可对角化。因为 \(100\) 为偶数，
\[
0^{100}=0,\qquad (-1)^{100}=1.
\]
于是 \(A^{100}\) 等于将 \(A\) 的特征值 \(0,-1\) 分别替换为 \(0,1\) 后得到的矩阵。直接由最小多项式也可看出
\[
A^{100}=A^2.
\]
计算
\[
A^2=
\begin{pmatrix}
-1&-4&2\\
0&1&0\\
-1&-2&2
\end{pmatrix}.
\]}
\answer{\[
A^{100}=
\begin{pmatrix}
-1&-4&2\\
0&1&0\\
-1&-2&2
\end{pmatrix}.
\]}
\end{solutionblock}

\begin{solutionblock}
\textbf{5.4 例 9}

\question{设 \(A=\begin{pmatrix}-2&-2&1\\2&x&-2\\0&0&-2\end{pmatrix}\)，\(B=\begin{pmatrix}2&1&0\\0&-1&0\\0&0&y\end{pmatrix}\)。求可逆矩阵 \(P\)，使 \(P^{-1}AP=B\)。}



\analysis{由 5.3 例 1 已知相似时
\[
x=3,\qquad y=-2.
\]
于是
\[
A=\begin{pmatrix}
-2&-2&1\\
2&3&-2\\
0&0&-2
\end{pmatrix},\qquad
B=\begin{pmatrix}
2&1&0\\
0&-1&0\\
0&0&-2
\end{pmatrix}.
\]
要求 \(P^{-1}AP=B\)，等价于
\[
AP=PB.
\]
若 \(P=(p_1,p_2,p_3)\)，则由 \(PB\) 的列关系得
\[
Ap_1=2p_1,\qquad Ap_2=p_1-p_2,\qquad Ap_3=-2p_3.
\]
可取
\[
p_1=(-1,2,0)^T,\qquad p_2=(1,0,0)^T,\qquad p_3=(-1,2,4)^T.
\]
于是
\[
P=\begin{pmatrix}
-1&1&-1\\
2&0&2\\
0&0&4
\end{pmatrix},
\qquad |P|=-8\ne0,
\]
并且直接验证
\[
P^{-1}AP=B.
\]}
\answer{可取
\[
P=\begin{pmatrix}
-1&1&-1\\
2&0&2\\
0&0&4
\end{pmatrix}.
\]}
\end{solutionblock}

\subsection{5.5 实对称矩阵相似对角化}

\begin{solutionblock}
\textbf{5.5 例 1}

\question{设二维向量 \(\alpha,\beta\) 满足 \(\|\alpha+\beta\|=\|\alpha-\beta\|\)，令 \(A=(\alpha,\beta)\)，判断 \(A\) 的列向量关系及 \(A^TA\) 的性质。}



\analysis{由
\[
\|\alpha+\beta\|=\|\alpha-\beta\|
\]
两边平方：
\[
(\alpha+\beta)^T(\alpha+\beta)
=(\alpha-\beta)^T(\alpha-\beta).
\]
展开得
\[
\alpha^T\alpha+2\alpha^T\beta+\beta^T\beta
=\alpha^T\alpha-2\alpha^T\beta+\beta^T\beta,
\]
故
\[
\alpha^T\beta=0.
\]
又 \(\alpha,\beta\) 均为单位列向量，所以矩阵
\[
A=(\alpha,\beta)
\]
的两列是标准正交向量，满足
\[
A^TA=E.
\]
因此 \(A\) 是正交矩阵。}
\answer{选 C。}
\end{solutionblock}

\begin{solutionblock}
\textbf{5.5 例 2}

\question{设 \(A\) 为四阶实对称矩阵，满足 \(A^2=-A\)，且 \(r(A)=3\)，判断 \(A\) 相似于哪个对角矩阵。}



\analysis{\(A\) 为四阶实对称矩阵，故正交相似于对角矩阵。设其特征值为 \(\lambda\)。由
\[
A^2=-A
\]
知
\[
\lambda^2=-\lambda,\qquad \lambda(\lambda+1)=0.
\]
所以特征值只能是 \(0\) 或 \(-1\)。又
\[
r(A)=3,
\]
说明非零特征值有 \(3\) 个，故 \(-1\) 出现三次，\(0\) 出现一次。}
\answer{选 D，即 \(A\sim \operatorname{diag}(-1,-1,-1,0)\)。}
\end{solutionblock}

\begin{solutionblock}
\textbf{5.5 例 3}

\question{设实对称矩阵 \(A\) 满足 \(A^3+A^2+A-3E=0\)，证明 \(A=E\)。}



\analysis{由于 \(A\) 为实对称矩阵，存在正交矩阵 \(Q\)，使
\[
Q^TAQ=\Lambda=\operatorname{diag}(\lambda_1,\lambda_2,\ldots,\lambda_n),
\]
其中 \(\lambda_i\) 均为实数。题设
\[
A^3+A^2+A-3E=0,
\]
所以每个特征值都满足
\[
\lambda^3+\lambda^2+\lambda-3=0.
\]
因式分解：
\[
\lambda^3+\lambda^2+\lambda-3
=(\lambda-1)(\lambda^2+2\lambda+3).
\]
二次因子 \(\lambda^2+2\lambda+3\) 的判别式为
\[
4-12=-8<0,
\]
没有实根。故实特征值只能是
\[
\lambda=1.
\]
于是 \(\Lambda=E\)，从而
\[
A=QEQ^T=E.
\]}
\answer{已证 \(A=E\)。}
\examnote{实对称矩阵的核心优势是“特征值全为实数、可正交对角化”，所以多项式方程可以转化为实数方程。}
\end{solutionblock}

\begin{solutionblock}
\textbf{5.5 例 4}

\question{设三阶实对称矩阵 \(A\) 的特征值为 \(-1,1,2\)。已知 \(\lambda_1=-1\) 的特征向量为 \(\alpha_1=(1,2,1)^T\)，\(\lambda_2=1\) 的特征向量为 \(\alpha_2=(-1,0,1)^T\)，求 \(\lambda_3=2\) 的特征向量。}



\analysis{实对称矩阵不同特征值对应的特征向量正交。已知
\[
\alpha_1=(1,2,1)^T,\qquad \alpha_2=(-1,0,1)^T
\]
分别属于 \(\lambda_1=-1,\lambda_2=1\)，所求 \(\lambda_3=2\) 的特征向量 \(\alpha_3\) 必须同时满足
\[
\alpha_3^T\alpha_1=0,\qquad \alpha_3^T\alpha_2=0.
\]
取叉乘：
\[
\alpha_1\times\alpha_2=(2,-2,2)^T.
\]
故可取
\[
\alpha_3=(1,-1,1)^T.
\]}
\answer{\[
\alpha_3=k(1,-1,1)^T,\qquad k\ne0.
\]}
\end{solutionblock}

\begin{solutionblock}
\textbf{5.5 例 5}

\question{设三阶实对称矩阵 \(A\) 每行元素之和均为 \(3\)，且 \(\lambda_1=\lambda_2=1\) 为二重特征值，求 \(A\) 属于 \(\lambda=1\) 的特征向量空间。}



\analysis{设 \(\mathbf{1}=(1,1,1)^T\)。每行元素之和均为 \(3\)，等价于
\[
A\mathbf{1}=3\mathbf{1}.
\]
因此 \(3\) 是 \(A\) 的一个特征值，对应特征向量为 \((1,1,1)^T\)。

题设还有二重特征值 \(\lambda_1=\lambda_2=1\)。由于 \(A\) 为实对称矩阵，不同特征值对应的特征向量正交，所以属于 \(\lambda=1\) 的特征向量必须与 \((1,1,1)^T\) 正交，即满足
\[
x_1+x_2+x_3=0.
\]
该平面的一组基为
\[
(1,-1,0)^T,\qquad (1,0,-1)^T.
\]}
\answer{\[
\lambda=3:\ \xi=k(1,1,1)^T,\ k\ne0;
\]
\[
\lambda=1:\ \xi=s(1,-1,0)^T+t(1,0,-1)^T,\ (s,t)\ne(0,0).
\]}
\end{solutionblock}

\begin{solutionblock}
\textbf{5.5 例 6}

\question{设 \(A=\begin{pmatrix}1&-2&2\\-2&-2&4\\2&4&-2\end{pmatrix}\)，求正交矩阵 \(Q\)，使 \(Q^{-1}AQ=Q^TAQ=\Lambda\)。}



\analysis{设
\[
A=\begin{pmatrix}
1&-2&2\\
-2&-2&4\\
2&4&-2
\end{pmatrix}.
\]
计算特征多项式：
\[
|\lambda E-A|=(\lambda-2)^2(\lambda+7).
\]
对应 \(\lambda=2\) 的特征子空间可取两个正交向量
\[
\xi_1=(-2,1,0)^T,\qquad \xi_2=(2,4,5)^T.
\]
对应 \(\lambda=-7\) 的特征向量可取
\[
\xi_3=(-1,-2,2)^T.
\]
归一化：
\[
q_1=\frac1{\sqrt5}(-2,1,0)^T,\qquad
q_2=\frac1{3\sqrt5}(2,4,5)^T,\qquad
q_3=\frac13(-1,-2,2)^T.
\]
令
\[
Q=(q_1,q_2,q_3).
\]
则 \(Q\) 为正交矩阵，且
\[
Q^TAQ=\operatorname{diag}(2,2,-7).
\]}
\answer{\[
Q=\begin{pmatrix}
-\frac2{\sqrt5}&\frac2{3\sqrt5}&-\frac13\\
\frac1{\sqrt5}&\frac4{3\sqrt5}&-\frac23\\
0&\frac5{3\sqrt5}&\frac23
\end{pmatrix},
\qquad
\Lambda=\operatorname{diag}(2,2,-7).
\]}
\end{solutionblock}

\begin{solutionblock}
\textbf{5.5 例 7}

\question{设 \(A\) 为三阶实对称矩阵，各行元素之和均为 \(3\)，且 \(\alpha_1=(1,-2,1)^T,\alpha_2=(0,1,-1)^T\) 是 \(Ax=0\) 的基础解系，求 \(A\) 及 \((A-\frac32E)^6\)。}



\analysis{设
\[
\alpha_1=(1,-2,1)^T,\qquad \alpha_2=(0,1,-1)^T.
\]
它们是 \(Ax=0\) 的两个线性无关解，所以
\[
A\alpha_1=0,\qquad A\alpha_2=0.
\]
又 \(A\) 为三阶实对称矩阵且各行元素之和均为 \(3\)，故
\[
A(1,1,1)^T=3(1,1,1)^T.
\]
注意
\[
\alpha_1^T(1,1,1)=0,\qquad \alpha_2^T(1,1,1)=0,
\]
故这三个向量构成一组正交特征向量。于是 \(A\) 在这组特征向量下的特征值为 \(0,0,3\)。这说明 \(A\) 是到 \((1,1,1)^T\) 方向的三倍投影：
\[
A=3\cdot \frac{(1,1,1)^T(1,1,1)}{(1,1,1)(1,1,1)^T}
=3\cdot\frac13
\begin{pmatrix}
1&1&1\\
1&1&1\\
1&1&1
\end{pmatrix}
=\begin{pmatrix}
1&1&1\\
1&1&1\\
1&1&1
\end{pmatrix}.
\]
矩阵 \(A-\frac32E\) 的特征值为
\[
-\frac32,\ -\frac32,\ \frac32.
\]
六次方后均变为
\[
\left(\frac32\right)^6=\frac{729}{64}.
\]
故
\[
\left(A-\frac32E\right)^6=\frac{729}{64}E.
\]}
\answer{\[
A=\begin{pmatrix}
1&1&1\\
1&1&1\\
1&1&1
\end{pmatrix},
\qquad
\left(A-\frac32E\right)^6=\frac{729}{64}E.
\]}
\end{solutionblock}
